\section{선별 문제 풀이 II}
\label{sec:polynomial-solutions-b}

\subsection*{문제 \thechapter.20: 다항식환의 중국인의 나머지 정리}

\begin{solution}
$\gcd(f,g)=1$이므로 베주 항등식에 의해 $af+bg=1$인 $a,b\in K[X]$가 존재한다. 따라서 $(f)+(g)=K[X]$이고 두 아이디얼은 서로소이다. 환의 중국인의 나머지 정리를 적용하면
\[
 K[X]/((f)\cap(g))
 \cong K[X]/(f)\times K[X]/(g).
\]
서로소인 주아이디얼에 대해
\[
 (f)\cap(g)=(fg)
\]
이므로 원하는 동형을 얻는다. 구체적인 준동형은
\[
 h+(fg)\longmapsto(h+(f),h+(g))
\]
이다.
\end{solution}

\subsection*{문제 \thechapter.21: 유리근과 인수분해}

\begin{solution}
유리근 $r/s$를 기약분수로 쓰면
\[
 r\mid6,\qquad s\mid2
\]
이므로 서로 다른 후보는
\[
 \pm1,\pm2,\pm3,\pm6,
 \quad
 \pm\frac12,\pm\frac32
\]
이다. 대입하면 $3/2$가 근이고
\[
 2X^3-3X^2+4X-6=(2X-3)(X^2+2).
\]
$X^2+2$는 $\Q$에 근이 없으므로 기약이다. 이것이 $\Q[X]$에서의 완전한 인수분해이다.
\end{solution}

\subsection*{문제 \thechapter.23: 아이젠슈타인 판정}

\begin{solution}
(a) $p=7$을 택한다. 최고차항계수 $1$은 $7$로 나누어지지 않고, $X^6$부터 상수항까지의 계수
\[
  0,0,0,14,0,21,7
\]
은 모두 $7$로 나누어진다. 상수항 $7$은 $49$로 나누어지지 않는다. 따라서 다항식은 $\Q[X]$에서 기약이다.

(b) $p=2$를 택한다. 최고차항계수 $3$은 짝수가 아니고, 나머지 계수
\[
  10,0,20,0,50
\]
은 모두 짝수이다. 상수항 $50$은 $4$로 나누어지지 않는다. 다항식의 내용은 $1$이므로 아이젠슈타인 판정법에 따라 $\Z[X]$와 $\Q[X]$에서 기약이다.
\end{solution}

\subsection*{문제 \thechapter.24: $X^4+1$의 기약성}

\begin{solution}
변수 이동은 기약성을 보존한다. 따라서
\[
 f(X)=X^4+1
\]
대신
\[
 f(X+1)=(X+1)^4+1
 =X^4+4X^3+6X^2+4X+2
\]
를 본다. 이 다항식은 $p=2$에 대해 아이젠슈타인 조건을 만족한다. 따라서 $f(X+1)$이 $\Q[X]$에서 기약이고, 역대입 $X\mapsto X-1$에 의해 $X^4+1$도 기약이다.
\end{solution}

\subsection*{문제 \thechapter.26: $\F_2$ 위의 낮은 차수 기약다항식}

\begin{solution}
일계수 기약다항식은 다음과 같다.
\[
\begin{array}{c|l}
\text{차수}&\text{기약다항식}\\
\hline
1&X,\ X+1\\
2&X^2+X+1\\
3&X^3+X+1,\ X^3+X^2+1\\
4&X^4+X+1,\ X^4+X^3+1,\\
 &X^4+X^3+X^2+X+1
\end{array}
\]
차수 $2,3$에서는 $0,1$이 근이 아닌지를 검사하면 된다.

차수 $4$의 일계수 기약다항식은 상수항이 $1$이고 $0,1$을 근으로 갖지 않아야 한다. 따라서
\[
  X^4+aX^3+bX^2+cX+1
\]
에서 $a+b+c=1$이어야 하며 후보는 네 개뿐이다. 한편 $\F_2[X]$의 유일한 일계수 기약 이차다항식은
\[
  q=X^2+X+1
\]
이다. 근이 없는 일계수 사차다항식이 가약이라면 두 기약 이차다항식의 곱이어야 하므로 반드시
\[
  q^2=X^4+X^2+1
\]
이다. 네 후보 중 이 다항식 하나를 제외하면 표에 적은 세 기약 사차다항식을 얻는다.
\end{solution}

\subsection*{문제 \thechapter.28: $K[X]/(X^2-1)$의 분해}

\begin{solution}
준동형
\[
 \Phi:K[X]\to K\times K,
 \qquad f\mapsto(f(1),f(-1))
\]
를 생각하자. 라그랑주 보간법 또는 직접 계산으로 $\Phi$는 전사이다. 실제로 $2$가 단위원이므로 임의의 $(a,b)$에 대해
\[
 f(X)=\frac{a+b}{2}+\frac{a-b}{2}X
\]
이면 $f(1)=a$, $f(-1)=b$이다.

핵은 $(X-1)\cap(X+1)$이다. 표수가 $2$가 아니므로 두 아이디얼은 서로소이고
\[
 \ker\Phi=(X-1)(X+1)=(X^2-1).
\]
제1동형정리에 의해
\[
 K[X]/(X^2-1)\cong K\times K.
\]
$X+(X^2-1)$은 $(1,-1)$에 대응한다.
\end{solution}

\subsection*{문제 \thechapter.30: $X^3-2$와 분해체}

\begin{solution}
(a) $X^3-2$는 $p=2$에 대해 아이젠슈타인 조건을 만족하므로 $\Q[X]$에서 기약이다.

(b) 따라서 $(X^3-2)$는 극대아이디얼이고
\[
 \Q[X]/(X^3-2)
\]
는 체이다.

(c) $\alpha=\sqrt[3]{2}$라 하면 이 몫체는 실수체 $\Q(\alpha)$와 동형이다. 그러나 복소수에서 세 근은
\[
 \alpha,\qquad \omega\alpha,\qquad \omega^2\alpha,
\]
여기서 $\omega=e^{2\pi i/3}$이고 $\omega^2+\omega+1=0$이다. $\Q(\alpha)$는 실수만 포함하므로 두 비실근을 포함하지 않는다.

(d) 원시 세제곱근 $\omega$를 추가하면
\[
 \Q(\alpha,\omega)
\]
가 세 근을 모두 포함한다. 이것이 $X^3-2$의 $\Q$ 위 분해체이다.
\end{solution}
